% 星族（综述）
% license CCBYSA3
% type Wiki

本文根据 CC-BY-SA 协议转载翻译自维基百科\href{https://en.wikipedia.org/wiki/Stellar_population}{相关文章}。

\begin{figure}[ht]
\centering
\includegraphics[width=6cm]{./figures/1ba3ca4c285278ba.png}
\caption{艺术家构想的银河系螺旋结构示意图，展示了巴德提出的恒星总体分类。螺旋臂中的蓝色区域由较年轻的Ⅰ族恒星组成，而中央核球中的黄色恒星则是较古老的Ⅱ族恒星。实际上，许多Ⅰ族恒星也与较老的Ⅱ族恒星混合分布在一起。} \label{fig_Stpo_1}
\end{figure}
1944 年，瓦尔特·巴德（Walter Baade）将银河系中的恒星分组划分为不同的恒星族群。在巴德文章的摘要中，他指出这种分类思想最早由扬·奥尔特（Jan Oort）于 1926 年提出。$^{[1]}$

巴德观察到，偏蓝色的恒星与螺旋臂有着很强的关联，而黄色恒星则主要分布在银河系中心核球附近以及球状星团之中。$^{[2]}$ 因此，他提出了两个主要划分：Ⅰ族恒星（Population I）和Ⅱ族恒星（Population II）；1978 年又增加了一个较新的、假设性的划分，即Ⅲ族恒星（Population III）。

在不同恒星族群之间，人们发现其各自观测到的恒星光谱存在显著差异。后来研究表明，这些差异极其重要，并且可能与恒星形成、观测到的运动学特性$^{[3]}$、恒星年龄，甚至螺旋星系和椭圆星系的演化有关。这三种简单的恒星族群在化学组成（即金属丰度）方面对恒星进行了有效划分。$^{[4][5][3]}$ 在天体物理学术语中，“金属”指的是所有比氦更重的元素，其中也包括诸如氧这样的化学非金属元素。$^{[6]}$

按定义而言，每一个恒星族群都表现出这样一种趋势：金属含量越低，恒星年龄越高。因此，宇宙中最早形成的恒星（金属含量极低）被归为Ⅲ族恒星，古老恒星（低金属丰度）被归为Ⅱ族恒星，而较新的恒星（高金属丰度）则被归为Ⅰ族恒星。$^{[7]}$ 太阳被认为是Ⅰ族恒星，是一颗较新的恒星，其金属丰度相对较高，约为 1.4\%。
\subsection{恒星演化}
对恒星光谱的观测表明，比太阳更古老的恒星所含的重元素数量少于太阳。$^{[3]}$ 这立即表明，金属丰度是在恒星世代更替过程中通过恒星核合成这一过程逐步演化而来的。
\subsubsection{第一代恒星的形成}
根据当前的宇宙学模型，在大爆炸中形成的全部物质主要由氢（75\%）和氦（25\%）组成，只有极其微小的一部分由锂、铍等其他轻元素构成。$^{[8]}$ 当宇宙冷却到足够低的温度时，第一代恒星作为Ⅲ族恒星（Population III）诞生，它们不含任何后期污染的重元素。人们推测，这一特性影响了它们的内部结构，使其恒星质量可达到太阳的数百倍。反过来，这些大质量恒星的演化也极为迅速，其核合成过程产生了最初的 26 种元素（即元素周期表中直到铁为止的元素）。$^{[9]}$

许多理论恒星模型表明，大多数高质量的Ⅲ族恒星会迅速耗尽燃料，并很可能以极其高能的对不稳定超新星（pair-instability supernovae）形式爆炸。这些爆炸会将其物质彻底散播开来，把金属抛射进星际介质（ISM），并被后续世代的恒星所吸收。它们的毁灭意味着，银河系中不应还能观测到高质量的Ⅲ族恒星。$^{[10]}$ 然而，在高红移星系中，可能仍能看到一些Ⅲ族恒星，其光起源于宇宙更早期的历史阶段。$^{[11]}$ 科学家已经发现了一颗极端贫金属恒星的证据，这颗恒星略小于太阳，位于银河系螺旋臂中的一个双星系统内。这一发现为观测更古老的恒星打开了可能性。$^{[12]}$

那些质量过大、无法产生对不稳定超新星的恒星，很可能会通过一种称为光致解体（photodisintegration）的过程坍缩成黑洞。在这一过程中，部分物质可能以相对论喷流的形式逃逸，这同样可能将最初的金属散布到宇宙之中。$^{[13][14][a]}$
\subsubsection{已观测恒星的形成}
迄今为止观测到的最古老恒星$^{[10]}$被称为Ⅱ族恒星（Population II），它们具有极低的金属丰度。$^{[16][7]}$ 随着后续世代恒星的诞生，新形成的恒星逐渐变得更加富含金属，因为它们形成所依赖的气体云吸收了由Ⅲ族恒星等前一代恒星制造并释放出的富金属尘埃。

当这些Ⅱ族恒星死亡时，它们通过行星状星云和超新星爆发将富含金属的物质回馈给星际介质，进一步丰富了星云，而新的恒星正是从这些星云中形成的。因此，最年轻的恒星，包括太阳在内，具有最高的金属含量，被称为Ⅰ族恒星（Population I）。
\subsection{瓦尔特·巴德的化学分类}
\subsubsection{Ⅰ族恒星}
\begin{figure}[ht]
\centering
\includegraphics[width=6cm]{./figures/401c0b6650ed6293.png}
\caption{第一族恒星——参宿七（Rigel）及其反射星云 IC 2118。} \label{fig_Stpo_2}
\end{figure}
Ⅰ族恒星（Population I stars）是三种恒星族群中最年轻、金属丰度最高的一类，更常见于银河系的旋臂之中。太阳被认为是一颗中等的Ⅰ族恒星，而类似太阳的 μ Arae 则富含更多的金属。$^{[17]}$（“富金属”一词用于描述金属丰度显著高于太阳、且这种差异不能用测量误差来解释的恒星。）

Ⅰ族恒星通常围绕银河中心做规则的椭圆轨道运动，其相对速度较低。早期曾假设，Ⅰ族恒星较高的金属丰度使其比另外两类恒星更有可能拥有行星系统，因为行星，尤其是类地行星，被认为是通过金属的吸积形成的。$^{[18]}$ 然而，对开普勒空间望远镜数据的观测发现，在具有不同金属丰度的恒星周围都存在小型行星，而只有较大的、可能的气态巨行星才更集中地分布在金属丰度相对较高的恒星周围——这一结果对气态巨行星形成理论具有重要影响。$^{[19]}$ 在中等Ⅰ族恒星与Ⅱ族恒星之间，还存在一个中间的盘族恒星群体。
\subsubsection{Ⅱ族恒星}
\begin{figure}[ht]
\centering
\includegraphics[width=6cm]{./figures/70684b47769d036f.png}
\caption{银河系。第二族恒星位于银河核球以及球状星团中。} \label{fig_Stpo_3}
\end{figure}
\begin{figure}[ht]
\centering
\includegraphics[width=8cm]{./figures/ef10f92a9e0305c0.png}
\caption{大爆炸后约一亿年时，第三族恒星群的艺术家印象图。} \label{fig_Stpo_4}
\end{figure}
Ⅱ族恒星（Population II），亦称贫金属恒星，是指其所含氦以外重元素相对较少的恒星。这些天体形成于宇宙更早的时期。中等Ⅱ族恒星常见于银河系中心附近的核球区域，而分布在银河晕中的Ⅱ族恒星则更为古老，因此其金属丰度也更低。球状星团中同样包含大量Ⅱ族恒星。$^{[20]}$

Ⅱ族恒星的一个显著特征是，尽管其总体金属丰度较低，但相对于铁（Fe），它们往往具有更高比例的α元素（即通过α过程产生的元素，如氧和氖），这一比例高于Ⅰ族恒星。现有理论认为，这是因为在这些恒星形成的时代，Ⅱ型超新星对星际介质的贡献更为重要，而Ⅰa型超新星所导致的金属富集发生在宇宙演化的较晚阶段。$^{[21]}$

科学家已在多个巡天项目中将这些最古老的恒星作为研究目标，包括 Timothy C. Beers 等人开展的 HK 物镜棱镜巡天 $^{[22]}$，以及 Norbert Christlieb 等人主持的 Hamburg–ESO 巡天 $^{[23]}$，后者最初是为搜寻暗弱类星体而启动的。迄今为止，这些项目已发现并详细研究了大约十颗极端贫金属恒星（UMP 恒星，如 Sneden 星、Cayrel 星、BD +17° 3248），以及目前已知最古老的三颗恒星：HE 0107-5240、HE 1327-2326 和 HE 1523-0901。Caffau 星在 2012 年利用斯隆数字巡天（SDSS）数据被发现时，被认定为当时已知金属丰度最低的恒星。然而，在 2014 年 2 月，借助 SkyMapper 天文巡天数据，人们又宣布发现了一颗金属丰度更低的恒星——SMSS J031300.36-670839.3。金属贫化程度稍低、但距离更近且更明亮、因而更早被发现的恒星还包括 HD 122563（红巨星）和 HD 140283（次巨星）。位于 ED-2 银河晕恒星流中的 Gaia BH3 是目前已知距离地球第二近的黑洞，其伴星同样是一颗极端贫金属恒星。鉴于该黑洞异常高的质量（约为 32.70 个太阳质量），这表明大质量Ⅱ族恒星在演化过程中经历的质量损失较少，从而能够形成质量超过由富金属前身形成的恒星黑洞质量上限的黑洞。$^{[24]}$
\subsubsection{Ⅲ族恒星}
\begin{figure}[ht]
\centering
\includegraphics[width=8cm]{./figures/8d03a94d9f092c45.png}
\caption{由美国国家航空航天局（NASA）的斯皮策空间望远镜成像的、可能来自第三族恒星的辉光} \label{fig_Stpo_5}
\end{figure}
Ⅲ族恒星（Population III stars）$^{[25]}$ 是一种假想中的恒星族群，其特征是质量极大、光度极高、温度极高，并且几乎不含任何“金属”，除非存在来自附近其他早期Ⅲ族超新星喷射物的少量混合。该术语最早由 Neville J. Woolf 于 1965 年提出。$^{[26][27]}$ 这类恒星很可能存在于宇宙的极早期阶段（即高红移时期），并可能开启了氢以外更重化学元素的产生过程，而这些元素是后来行星形成以及我们所知生命出现所必需的。$^{[28][29]}$

Ⅲ族恒星的存在主要是根据物理宇宙学推断出来的，但迄今尚未被直接观测到。在宇宙极其遥远区域的一些引力透镜星系中，人们已经发现了支持其存在的间接证据。$^{[30]}$ 它们的存在或可解释一个事实：在类星体的发射光谱中观测到了重元素，而这些元素不可能在大爆炸中形成。$^{[9]}$ 人们还认为它们可能是暗弱蓝星系的组成部分。这些恒星很可能触发了宇宙的再电离时期，即构成大部分星际介质的氢气发生的一次重要相变。对星系 UDFy-38135539 的观测表明，它可能在这一再电离过程中发挥了作用。欧洲南方天文台在一个极为明亮的星系 Cosmos Redshift 7 中发现了一处早期恒星族群的高亮区域，该星系处于大爆炸后约 8 亿年的再电离时期，红移为 z = 6.60；而该星系的其余部分则包含一些较晚、颜色更红的Ⅱ族恒星。$^{[28][31]}$ 一些理论认为，Ⅲ族恒星可能存在两个世代。$^{[32]}$
\begin{figure}[ht]
\centering
\includegraphics[width=8cm]{./figures/45c01bd28c2f4623.png}
\caption{大爆炸后约 4 亿年时第一代恒星的艺术想象图} \label{fig_Stpo_6}
\end{figure}
当前理论在“第一代恒星是否具有极高质量”这一问题上仍存在分歧。一种观点认为，这些恒星的质量远大于当今恒星，可能达到数百个太阳质量，甚至高达 1,000 个太阳质量。这样的恒星寿命极短，仅可存在约 200 万至 500 万年。$^{[33]}$ 由于缺乏重元素以及大爆炸后星际介质温度更高，这种超大质量恒星被认为在当时是可能形成的。$^{[citation\ needed]}$ 相反，2009 年和 2011 年提出的另一类理论认为，第一代恒星群可能由一颗大质量恒星及其周围的若干颗较小恒星组成。$^{[34][35][36]}$ 这些较小的恒星如果仍留在其诞生星团中，会继续吸积气体而无法存活至今；但 2017 年的一项研究指出，如果一颗质量为 0.8 个太阳质量（$M_{\odot}$）或更低的恒星在尚未进一步增质量之前就被抛射出其诞生星团，那么它可能一直存活到今天，甚至可能仍存在于我们的银河系中。$^{[37]}$

对极低金属丰度的Ⅱ族恒星（如 HE 0107-5240）的数据分析表明，这类恒星所含金属很可能由Ⅲ族恒星产生，这暗示这些无金属恒星的初始质量约为 20–130 个太阳质量。$^{[38]}$ 另一方面，对与椭圆星系相关的球状星团的分析则表明，其金属组成可能源自成对不稳定性超新星，而这类超新星通常与极大质量恒星相关。$^{[39]}$ 这也解释了为何尽管针对较小质量Ⅲ族恒星的模型已被提出，但迄今尚未观测到零金属的低质量恒星。$^{[40][41]}$ 有学者提出，包含零金属红矮星或褐矮星（可能由成对不稳定性超新星产生$^{[16]}$）的星团或可作为暗物质候选体。$^{[42][43]}$ 然而，通过引力微透镜搜索此类 MACHO 的尝试迄今未获得正面结果。$^{[citation\ needed]}$

Ⅲ族恒星被认为是早期宇宙中黑洞的“种子”。与直接坍缩黑洞等高质量黑洞种子不同，它们最初形成的应是低质量黑洞。若这些黑洞能够增长至超出预期的质量，它们可能演化为“准恒星”，即另一类假想中的大质量黑洞种子，被认为存在于宇宙早期、氢和氦尚未被重元素污染之前。

探测Ⅲ族恒星是 NASA 詹姆斯·韦布空间望远镜的重要科学目标之一。$^{[44]}$

2022 年 12 月 8 日，天文学家报告称，在一个名为 RX J2129–z8He II 的高红移星系中，可能探测到了Ⅲ族恒星的迹象。$^{[45][46]}$

2025 年 10 月 27 日，天文学家在《天体物理学报通讯》（Astrophysical Journal Letters）上发表论文，报告称 LAP1-B 可能是首个被观测到的Ⅲ族恒星，其红移为 z = 6.6，并且似乎满足将恒星判定为Ⅲ族恒星的三个最关键判据。$^{[47][48]}$
\subsection{另见}
\begin{itemize}
\item 天文对象列表
\item 恒星列表
\item 躲猫猫星系（Peekaboo Galaxy）
\end{itemize}

\subsection{注释}
a.有人提出，近期的超新星 SN 2006gy 和 SN 2007bi 可能属于对不稳定超新星（pair-instability supernovae），即由此类超大质量的第三族恒星发生爆炸所致。Clark（2010）推测，这些恒星可能在相对较近的宇宙时期于矮星系中形成，因为这些星系主要包含原初的、无金属的星际物质。过去发生在这些小星系中的超新星，可能以足够高的速度将富含金属的物质抛射出去，使其逃逸出星系，从而使这些小星系的金属丰度长期保持在极低水平。$^{[15]}$
